\Needspace{10\baselineskip}
\originalchapter{麦克斯韦方程与电磁势}
\originalsection{电荷连续性与位移电流}
电荷局部守恒是 $\partial_t\rho+\nabla\cdot\jj=0$. 单独使用稳恒安培定律 $\nabla\times\bb=\mu_0\jj$ 并取散度会错误地要求 $\nabla\cdot\jj=0$, 因而它不能描述充电电容器. 加入位移电流后, 真空方程为
\begin{align}
\nabla\cdot\ee&=\rho/\eps,& \nabla\cdot\bb&=0,\\
\nabla\times\ee&=-\partial_t\bb,&
\nabla\times\bb&=\mu_0\jj+c^{-2}\partial_t\ee.
\end{align}
对最后一式取散度, 第一式使它恰好成为连续性方程. 位移电流不是额外流动的带电粒子, 而是时变电场对磁场环流的贡献.

在运动的导线回路上, 单位电荷受到 $\ee+\vct u\times\bb$, 所以电动势为 $\mathcal E=\oint(\ee+\vct u\times\bb)\cdot\dd\vct l$. 用曲面输运定理可合成 $\mathcal E=-\dd\Phi_B/\dd t$. 只把 $\partial_t\bb$ 积分会漏掉导线运动造成的面积变化.

\originalsection{规范自由与波动方程}
由无磁单极方程取 $\bb=\nabla\times\vct A$. 法拉第定律随后给 $\ee=-\nabla V-\partial_t\vct A$. 势并非唯一: 对任意充分光滑的 $\Lambda$, $\vct A'=\vct A+\nabla\Lambda$, $V'=V-\partial_t\Lambda$, 场保持不变. 规范自由改变描述, 不改变可测量场.

解耦可以逐式看出. 记$\mathcal C=\nabla\cdot\vct A+c^{-2}\partial_tV$, $L=\nabla^2-c^{-2}\partial_t^2$. 高斯方程代入势给$-\nabla^2V-\partial_t(\nabla\cdot\vct A)=\rho/\eps$. 安培方程用$\nabla\times\nabla\times\vct A=\nabla(\nabla\cdot\vct A)-\nabla^2\vct A$给
\[
\nabla(\nabla\cdot\vct A)-\nabla^2\vct A
=\mu_0\jj-c^{-2}\nabla\partial_tV-c^{-2}\partial_t^2\vct A.
\]
分别整理为$LV=-\rho/\eps-\partial_t\mathcal C$, $L\vct A=-\mu_0\jj+\nabla\mathcal C$. 因而选择$\mathcal C=0$恰好消去两个耦合源, 而不是无缘由地抛弃一个项.

从任意势达到洛伦兹规范, 需要选择满足
\[
(\nabla^2-c^{-2}\partial_t^2)\Lambda=-\nabla\cdot\vct A-c^{-2}\partial_tV
\]
的规范函数, 并选择与问题相容的初始及边界条件. 达到后仍可添加满足齐次波动方程的 $\Lambda$, 这是残余规范自由. 库仑规范在动态问题中给 $\nabla^2V=-\rho/\eps$, 但矢量势还含横向电流的波动传播; 不能把第6章的稳恒泊松方程直接用于所有动态问题.

洛伦兹规范 $\nabla\cdot\vct A+c^{-2}\partial_tV=0$ 使源方程解耦为
\[
(\nabla^2-c^{-2}\partial_t^2)V=-\rho/\eps,\qquad
(\nabla^2-c^{-2}\partial_t^2)\vct A=-\mu_0\jj.
\]
选择无入射自由场的推迟解,
\[
V(\rr,t)=\frac1{4\pi\eps}\int\frac{\rho(\rr',t-R/c)}R\dd^3x',\quad
\vct A(\rr,t)=\frac{\mu_0}{4\pi}\int\frac{\jj(\rr',t-R/c)}R\dd^3x',
\quad R=|\rr-\rr'|.
\]
波动算子的推迟格林函数核为 $\delta(t-t'-R/c)/(4\pi R)$. 本稿记 $L=\nabla^2-c^{-2}\partial_t^2$, 此核满足 $LG=-\delta^3(\rr-\rr\,\!\prime)\delta(t-t\prime)$. 时间积分给出上式; 静态极限恢复库仑势. 齐次波动解仍可叠加, 所以仅写麦克斯韦方程还不等于指定辐射边界条件.

\originalsection{点电荷的推迟势}
假定粒子速度始终小于光速, 观察点不在粒子世界线上; 未正则化的点源自场不能直接用于求自力. 粒子位置为 $\rr_p(t)$, 推迟时间 $t_r$ 满足 $t-t_r=|\rr-\rr_p(t_r)|/c$. 定义 $\vct R=\rr-\rr_p(t_r)$、$R=|\vct R|$、$\nn=\vct R/R$、$\vct\beta=\vct v(t_r)/c$, 以及 $\kappa=1-\nn\cdot\vct\beta$. 插入点源并积分时间狄拉克分布时, 根的导数产生 $1/\kappa$, 因而 Li\'enard--Wiechert 势为
\[
V=\frac q{4\pi\eps}\left[\frac1{\kappa R}\right]_{t_r},\qquad
\vct A=\frac{\vct v(t_r)}{c^2}V.
\]
忽略该雅可比因子会错误地把移动点源当成“在旧位置静止”的库仑电荷. 推迟时间对观测时间及位置也有依赖, 微分势时要用 $\partial_t t_r=1/\kappa$, $\nabla t_r=-\nn/(c\kappa)$.

相应电场为
\begin{align*}
\ee=\frac q{4\pi\eps}\left[
\frac{(1-\beta^2)(\nn-\vct\beta)}{\kappa^3R^2}
+\frac{\nn\times[(\nn-\vct\beta)\times\dot{\vct\beta}]}{c\kappa^3R}
\right]_{t_r},\qquad
\bb=\frac1c\nn\times\ee.
\end{align*}
第一项是速度场, 随 $R^{-2}$ 衰减; 第二项是加速度场, 随 $R^{-1}$ 衰减并可把有限能流送到无穷远. 下文在中文册内从推迟核的雅可比开始, 完整展开势的微分; 辐射部分的非相对论推导在第\ref{ch:radiation}章另行完成.

\originalsection{从推迟核到移动点源的雅可比}
先把四维推迟积分中的时间变量保留, 点源的空间积分完成后, 标势成为
$V=q(4\pi\eps)^{-1}\int\delta[t-t'-R(t')/c]\dd t'/R(t')$.
对 $h(t')=t-t'-R(t')/c$, 有 $\dd R/\dd t'=-\nn\cdot\vct v$, 所以在根 $t_r$ 处 $h'=-1+\nn\cdot\vct\beta=-\kappa$. 亚光速保证 $\kappa>0$, 因而简单根分布公式给 $\delta(h)=\delta(t'-t_r)/\kappa$. 这产生 $V=q/(4\pi\eps\kappa R)$; 矢量势以相同雅可比给 $\vct A=\vct\beta V/c$. 观察点固定时 $h$ 对时间微分, 或观察时间固定时对位置微分, 分别得到前文的 $\partial_tt_r=1/\kappa$ 与 $\nabla t_r=-\nn/(c\kappa)$.

\originalsection{实际微分 Li\'enard--Wiechert 势}
为让速度和加速度项的分母能够逐步核对, 记 $F=\kappa R=R-\vct\beta\cdot\vct R$, $K_q=q/(4\pi\eps)$. 下文 $\dot{\vct\beta}$ 指对源时间 $t_r$ 求导, 所有量在 $t_r$ 取值. 推迟位置的导数为
\[
 \partial_t\vct R=-\frac{c\vct\beta}{\kappa},\qquad
 \partial_j R_i=\delta_{ij}+\frac{\beta_in_j}{\kappa}.
\]
对 $F=R-\vct\beta\cdot\vct R$ 用乘积法则, 注意 $\partial_t\vct\beta=\dot{\vct\beta}/\kappa$、$\partial_j\beta_i=-\dot\beta_in_j/(c\kappa)$, 得
\begin{align*}
 \partial_tF&=\frac{c(\beta^2-\nn\cdot\vct\beta)-R\nn\cdot\dot{\vct\beta}}{\kappa},\\
 \nabla F&=-\vct\beta+\frac{1-\beta^2+(R/c)\nn\cdot\dot{\vct\beta}}{\kappa}\nn.
\end{align*}
这两行同时包含 $R$ 改变、源速度改变和推迟时间改变, 任何一项都不能当成常量删掉. 由 $V=K_q/F$、$\vct A=K_q\vct\beta/(cF)$,
\[
 \ee=K_q\left[\frac{\nabla F+\vct\beta\partial_tF/c}{F^2}
 -\frac{\dot{\vct\beta}}{c\kappa F}\right].
\]
通分后, 分子中的速度部分是 $(1-\beta^2)(\nn-\vct\beta)$, 加速度部分是
$(R/c)[(\nn-\vct\beta)(\nn\cdot\dot{\vct\beta})-\kappa\dot{\vct\beta}]$,
共同分母为 $\kappa^3R^2$. 三重积恒等式将方括号写成 $\nn\times[(\nn-\vct\beta)\times\dot{\vct\beta}]$, 恰得前文的完整场式.

磁场同样直接取旋度. $\nabla\times\vct\beta=\nabla t_r\times\dot{\vct\beta}=-\nn\times\dot{\vct\beta}/(c\kappa)$, 故
$\bb=-K_q\nabla F\times\vct\beta/(cF^2)-K_q\nn\times\dot{\vct\beta}/(c^2\kappa F)$.
若把 $\nabla F=-\vct\beta+\lambda\nn$ 中的系数记 $\lambda$, 上两导数满足 $\lambda=1-\partial_tF/c$. 代回即给 $\bb=\nn\times\ee/c$. 所以这条关系对该点电荷的全场成立, 并非仅从平面波类推.

对静止电荷 $\vct\beta=\dot{\vct\beta}=0$, 速度场恢复库仑场. 对匀速电荷只有 $1/R^2$ 项; 对非相对论加速电荷, $\kappa\simeq1$、$\dot{\vct\beta}=\vct a/c$, 加速度场为 $q\nn\times(\nn\times\vct a)/(4\pi\eps c^2R)$. 它的场幅随 $R^{-1}$ 衰减, 乘球面积后仍有有限功率, 是辐射的几何原因. 本段的中文求导补足原第15讲的数学内容, 不需要读者回英文原件补步骤.


本章依据 8.07 第3讲、作业9和第15讲数学部分. 第15讲原件含明确许可例外的 YouTube 图片, 故未把整份原件打包, 本稿不使用这些图片.
